\chapter{安全、水印与隐私}

\paragraph*{本章阅读锚点}
安全方向的 3DGS 方法要从两个相反的问题读起。一方面，3DGS 正在变成可交易、可下载、可二次渲染的三维资产，显式高斯参数比 NeRF 权重更容易被复制、篡改和分析，因此需要水印、隐写和版权验证；另一方面，3DGS 的训练和渲染又暴露出新的攻击面，例如自适应 densification 会被投毒图像推高计算成本，目标视角可以被植入幻觉物体，重建后的 3D 资产也可能继承 2D 对抗噪声。本章的锚点不是原始 3DGS 的渲染质量，而是它的参数结构、训练流程和资产分发方式：每篇方法都应看它把信息或攻击接在了哪里，是球谐颜色、锚点特征、渲染图像、输入训练图像、密度控制，还是物理重光照模拟。

\section{对抗攻击与防御}

对抗方向首先揭示的是 3DGS 与 NeRF 不同的安全表面。NeRF 攻击多发生在网络权重或体渲染输出上，而 3DGS 有显式高斯、可见性排序、densification 和快速可微渲染；这些机制既能被利用来植入目标视角幻觉，也能被扰动触发额外分裂和克隆，造成计算拒绝服务。本节按攻击目标组织：先是视角/场景混淆，再是计算成本攻击，再到物理对抗伪装和对应防御。



% METHOD: 2410.08190
\paragraph{Poison-splat（计算成本攻击）（2024.10）}
Poison-splat~\cite{2410_08190} 抓住的是原始 3DGS 的 densification 机制。3DGS 会根据重建误差、梯度和可见性不断克隆、分裂或删除高斯，这本来是提高质量的关键；Poison-splat 的改变是把这一机制变成攻击面，通过微小投毒扰动让训练过程误以为场景中存在大量欠重建细节，从而不断生成更多高斯，推高显存和训练时间，严重时造成服务端拒绝服务。

具体做法形式化为双层优化：内层是受害者按普通流程训练 3DGS，外层优化输入图像扰动，使最终高斯数量、训练成本或内存占用最大化。由于完整双层优化代价过高，方法用三个近似降低攻击成本：用高斯数量增长率等代理指标近似真实计算成本，用代理渲染模型快速估计扰动效果，并在需要隐蔽时加入扰动幅度约束。项目 README 标明其 ICLR spotlight 定位和官方实现；项目实现围绕原始 3DGS 依赖和攻击组织，而不是提出新的渲染表示。它相对普通图像对抗攻击的关键差异在于，攻击成功的判据不是分类错误，而是训练系统被推向最坏复杂度。后续 RemedyGS 正是沿着这个威胁模型提出防御。
\textbf{实现细节：} Poison-splat（2024.10） 修改了 3DGS 的训练数据输入环节——用经 TV 最大化优化的投毒图像替代原始干净图像——3DGS 的自适应密度控制机制保持不变，但被恶意利用为攻击的放大器。该方法揭示了 3DGS 一个此前未被关注的安全盲区：攻击目标不是重建质量，而是计算资源，极端情况下可导致拒绝服务（DoS，Denial of Service）。
\textbf{方法信息：} 代码：\url{https://github.com/jiahaolu97/poison-splat}。
\textbf{设计动机与边界：} 相对原始 3DGS 以独立高斯和图像重建误差为主的设定，Poison-splat 把问题落在以下可实现的改动上：Poison-splat利用3DGS自适应密度控制策略的灵活性作为攻击向量。攻击目标为最大化3DGS训练中的高斯数量（即GPU内存和计算时间）。形式化为双层优化问题：外层优化投毒图像扰动，内层优化3DGS训练过程。三个关键策略：(1)攻击目标近似：用代理指标（如高斯数量增长率）近似实际攻击目标；(2)代理模型渲染：使用简化模型快速评估攻击效果，避免完整3DGS训练的巨大开销；(3)可选约束优化：限制扰动幅度保持攻击隐蔽性。投毒图像在视觉上与原始图像接近，但触发3DGS在欠重建区域大量分裂和克隆高斯。其设计目的不是单独增加一个模块，而是用上述表示、监督或渲染约束消除原流程在目标任务中的具体失配。首次揭示3DGS训练的计算成本安全漏洞，提出Poison-splat攻击方法通过数据投毒将3DGS训练推向最坏计算复杂度，极端情况下可导致拒绝服务攻击适用时还应保留论文给出的边界：攻击效果受限于3DGS的自适应密度控制阈值设置；高扰动幅度可能被检测到
% METHOD: 2412.02803
\paragraph{Gaussian Splatting Under Attack（从 2D 对抗噪声到 3DGS 资产）（2024.12）}
Gaussian Splatting Under Attack~\cite{2412_02803} 研究的是另一类攻击入口：不直接修改 3DGS 训练机制，而是在重建前对输入图像施加对抗噪声，再观察这种噪声是否会被 3DGS 固化到三维资产中。它的锚点是 IFGSM 图像攻击和 CLIP 视觉语言分类；改变在于加入 SAM 掩码，把扰动限制在目标物体区域，然后用这些被扰动的多视图图像训练 3DGS。

具体流程是三步：先用 SAM 为每张输入图像得到目标物体掩码；再在 IFGSM 的迭代梯度更新中乘上掩码，只让噪声落在物体区域，背景尽量不被破坏；最后用扰动后的图像重建 3DGS，并把渲染图再次送入 CLIP 分类。这样做回答了一个很实际的问题：如果数据采集阶段被攻击，重建出的 3DGS 是否会把攻击带到未见视角。实验在 CO3D 物体上显示，训练视角和测试视角的 CLIP Top-1 都明显下降，说明对抗扰动可以跨过重建流程进入 3D 表示。它不像 GaussTrap 那样定义目标视角后门，也不像 Poison-splat 那样追求计算成本，而是把 3DGS 作为 2D 对抗噪声的“传播介质”来评估安全风险。
\textbf{实现细节：} 前面的攻击方法分别从训练数据操纵（CLOAK）、参数空间扰动（DAGGER）、物理对抗物体（3DGAA/PGA/R-PGA）和前馈模型输入（AdvSplat）等角度攻击 3DGS 管线，而另一条攻击路线瞄准的是 3DGS 的更上游——从源图像到 3D 模型的完整生成管线。与上述方法（攻击 3DGS 场景本身、其训练数据或前馈模型输入）不同，M-IFGSM（2024.12） 攻击的是 3D 生成管线的输入端——在源图像上生成对抗噪声，使基于 CLIP（Contrastive Language-Image Pre-training） 视觉语言模型的 3D 生成系统产出可逃避检测的 3D 物体。攻击方法至此已经覆盖了从源图像到最终输出的完整攻击链。接下来转向防御方法的研究。
\textbf{实现补充：} 实现上还要核对输入、输出、参数、特征、初始化、训练损失、正则、渲染或推理接口，以及优化中的剪枝、分裂、采样、编码或解码步骤；这些抓手决定该方法到底改了 3DGS 管线的哪一环。
\textbf{方法信息：} Amazon 参与了这篇论文。
\textbf{设计动机与边界：} 相对原始 3DGS 以独立高斯和图像重建误差为主的设定，Gaussian Splatting Under Attack 把问题落在以下可实现的改动上：M-IFGSM方法包含三个阶段。(1)掩码生成：使用SAM（Segment Anything Model）从输入图像中提取目标物体的分割掩码，精确标定攻击区域。(2)对抗扰动生成：在传统IFGSM基础上引入掩码约束，仅对掩码区域施加对抗扰动，保持背景完整性。支持定向和非定向攻击，通过交叉熵损失计算梯度，迭代更新对抗噪声，设置损失阈值20.0作为早停条件。扰动范围限制在RGB值[0,255]内。(3)3D模型重建：从对抗图像使用3D Gaussian Splatting重建3D模型，对抗噪声被嵌入到3D表示中。攻击对训练和测试视角的渲染图像均有效。其设计目的不是单独增加一个模块，而是用上述表示、监督或渲染约束消除原流程在目标任务中的具体失配。首次系统研究3D高斯泼溅模型的对抗攻击安全性。提出M-IFGSM方法，针对CLIP视觉语言模型生成聚焦于目标物体区域的对抗噪声，在人类几乎不可察觉的情况下将3D模型渲染的CLIP Top-1准确率从95.4%降至12.5%（训练图像），揭示了3D视觉模型在安全关键应用中的严重脆弱性。适用时还应保留论文给出的边界：攻击需要目标物体的分割掩码；仅评估了CLIP模型；未提出防御机制；实验仅限于CO3D数据集的8个类别。
% METHOD: 2504.20829
\paragraph{GaussTrap（目标视角场景混淆）（2025.04）}
GaussTrap~\cite{2504_20829} 的锚点是原始 3DGS 的逐视角重建损失。普通训练会让所有训练视角共同约束同一组高斯，因此单独向一个视角塞入恶意内容容易破坏多视图一致性；GaussTrap 的改动是把训练过程拆成攻击、稳定化和正常训练三个阶段，让恶意内容只在指定攻击视角出现，同时让相邻和正常视角继续维持高质量重建。

具体做法是先定义三类相机：攻击视角用于呈现预设恶意图像，稳定化视角由攻击视角附近的小偏移相机组成，正常视角来自原始训练集。攻击阶段把目标视角的渲染结果拉向恶意内容；稳定化阶段使用 Viewpoint Ensemble Stabilization，用多个邻近视角缓冲单视角注入造成的局部几何和颜色不连续；正常阶段再用普通 L1 与 D-SSIM 组合损失恢复非攻击视角的渲染质量。它相对传统 2D 后门攻击的差异在于，攻击对象不是单张图像分类器，而是一个可从任意视角渲染的高斯场，因此必须同时优化“目标角度有效”和“其它角度隐蔽”这两个目标。论文在 Blender 与 Mip-NeRF360 场景上验证了这种三阶段调度能把恶意视图嵌入模型，同时让训练视角和测试视角保持较高 PSNR/SSIM。
\textbf{实现细节：} 与上述在训练数据层面实施攻击不同，后门攻击直接在模型训练过程中植入恶意行为，属于更隐蔽的威胁类别。Poison-splat 投毒训练数据攻击计算资源，PatchPoison 和 PoInit-of-View 攻击 SfM 特征匹配破坏重建质量，但这些方法的攻击效果在训练完成后即固定——它们影响的是模型的被动属性（计算开销或重建精度）。GaussTrap（2025.04） 首次系统研究 3DGS 后门威胁，将攻击范式从上述"一次性投毒"转向"按需激活"的后门攻击——在模型训练过程中植入仅在特定视角触发的恶意渲染行为。攻击目标不再是全局的计算资源或重建质量，而是推理阶段特定视角的场景语义。
\textbf{实现补充：} 实现上还要核对输入、输出、参数、特征、初始化、训练损失、正则、渲染或推理接口，以及优化中的剪枝、分裂、采样、编码或解码步骤；这些抓手决定该方法到底改了 3DGS 管线的哪一环。
\textbf{方法信息：} Huawei 参与了这篇论文。
\textbf{设计动机与边界：} 相对原始 3DGS 以独立高斯和图像重建误差为主的设定，GaussTrap 把问题落在以下可实现的改动上：GaussTrap是一种针对3D Gaussian Splatting模型的新型投毒攻击方法。方法包含三阶段优化管线：1）攻击阶段：在目标攻击视角嵌入恶意视图，使模型在特定角度展现预设恶意行为；2）稳定化阶段：引入视角扩展稳定化（VES）机制，利用多个微小偏移的相机视角来缓解单一恶意视角引起的不稳定，保持场景重建的连续性和一致性；3）正常训练阶段：确保模型在非攻击视角上生成高质量渲染结果，平衡攻击危害性与隐蔽性。每个阶段采用T次迭代，随机选择目标相机视角，通过高斯分布G和可微渲染器R渲染图像，使用L1损失和D-SSIM结构相似性损失的加权组合优化模型参数。方法在合成和真实世界数据集上验证了其鲁棒性和实用性。其设计目的不是单独增加一个模块，而是用上述表示、监督或渲染约束消除原流程在目标任务中的具体失配。首次系统性研究3DGS流水线中的后门威胁。提出GaussTrap投毒攻击方法，通过三阶段优化管线（攻击、稳定化、正常训练）在特定视角注入恶意场景混淆，同时保持非目标视角的高质量渲染。引入视角扩展稳定化（VES）机制缓解单视角恶意注入导致的不稳定性，在攻击有效性与隐蔽性之间取得平衡。适用时还应保留论文给出的边界：仅考虑了特定类型的后门攻击，对其他攻击形式（如全局投毒）的防御未涉及；防御策略未讨论；仅针对3DGS框架，对NeRF等其他表示的适用性未验证。
% METHOD: 2510.02314
\paragraph{StealthAttack（密度引导幻觉注入）（2025.10）}
StealthAttack~\cite{2510_02314} 可以看成 GaussTrap 之后更直接利用 3DGS 显式点云结构的投毒路线。相对只在图像或训练视角上加入后门的基线，它把攻击锚点放到高斯分布本身：目标仍是在目标视角渲染出幻觉物体，并保持非目标视角接近原模型；但它不只调度视角损失，而是先分析高斯点云密度，寻找低密度、低约束区域，把攻击内容放到更不容易被其它视角发现的位置。

具体做法有两部分。第一部分是密度引导点云攻击：对高斯中心做核密度估计，找出从目标视角可见、但在其它视角约束较弱的低密度区域，再向这些区域注入或调整高斯，使幻觉物体在目标视角形成清晰外观。第二部分是视图一致性破坏：对训练图像加入自适应噪声，使 3DGS 的多视图一致性约束变松，从而给目标视角幻觉留下优化空间。训练目标同时包含目标视角与幻觉图像的接近项，以及非目标视角与原始渲染的保持项；输出则是一个在特定视角触发幻觉、其它视角尽量正常的高斯模型。它还提出基于 KDE 的评估协议，用密度和视角重叠衡量攻击难度。与直接把幻觉贴到训练图像相比，这种方法更符合 3DGS 的显式表示特点：它不是让所有视角都学到同一个假物体，而是利用高斯分布中的稀疏区域让假物体只在特定观察条件下成立。
\textbf{实现细节：} 后门攻击通过三阶段管线植入恶意视角，而另一种更具隐蔽性的策略则是利用 3DGS 点云本身的稀疏性来隐藏攻击载荷。与 GaussTrap 需要从零开始三阶段训练不同，StealthAttack（2025.10） 直接操纵已训练好的 3DGS 点云，无需重新训练。这意味着攻击者可以在模型部署后随时注入恶意内容，威胁时间线从"训练期"延伸至"全生命周期"。StealthAttack 的形式化攻击目标为：给定原始 3DGS 点云 $G$，寻找投毒后的点云 $G^{*}$，使得投毒视角 $v_p$ 的渲染结果 $\mathcal{R}(G^{*}, v_p)$ 包含攻击者指定的幻觉物体 $O_*$（即 $\|\mathcal{R}(G^{*}, v_p) - I_*\|_2^{}2$ 最小化，其中 $I_*$ 为 $G$ 在 $v_p$ 的原始渲染与 $O_*$ 的叠加），同时所有非投毒视角 $v_k \neq v_p$ 的渲染与干净模型保持一致 $\|\mathcal{R}(G^{*}, v_k) - \mathcal{R}(G, v_k)\|_2^{}2 \approx 0$。StealthAttack 的设计动机源于对两种朴素方案的失败分析：（1）直接在目标视角的训练图像上叠加幻觉物体并重新训练——3DGS 的多视角一致性机制会将其视为噪声（违反一致性约束）并自动消除；（2）将幻觉物体直接反投影到高斯点云中——难以确定正确深度，可能被现有几何遮挡。由此 StealthAttack 得出关键结论：有效的 3DGS 投毒必须同时利用点云的显式表示特性和稀疏性结构，将幻觉物体放置在正常视角不可见的 3D 区域。从后门攻击和点云操纵转向下游检测器攻击，攻击者的目标从破坏重建质量扩展到欺骗感知系统。为此，StealthAttack 提出密度引导的点云攻击策略，包含场景空间分析、连续密度估计和最优位置选择三个步骤。
\textbf{实现补充：} 实现上还要核对输入、输出、参数、特征、初始化、训练损失、正则、渲染或推理接口，以及优化中的剪枝、分裂、采样、编码或解码步骤；这些抓手决定该方法到底改了 3DGS 管线的哪一环。
\textbf{方法信息：} NVIDIA 参与了这篇论文。



% METHOD: 2511.22147
\paragraph{RemedyGS（计算成本攻击防御）（2025.11）}
RemedyGS~\cite{2511_22147} 是对 Poison-splat 威胁模型的直接回应。若攻击来自用户上传的多视图图像，最保守的防御可以是限制高斯数量或对所有图像平滑，但前者会伤害正常重建质量，后者会把正常用户的细节也抹掉。RemedyGS 的改变是把防御前置到 3DGS 训练之前，并采用选择性处理：先检测哪些输入图像带有投毒纹理，再只净化这些图像，然后把净化后的图像交给未修改的 3DGS 训练。

具体做法包含检测器和净化器。检测器利用 Poison-splat 类攻击常引入异常高频纹理这一现象，用轻量卷积网络做干净/投毒二分类；被判定为干净的图像直接通过，被判定为投毒的图像进入净化器。净化器采用带跳连的编码器-解码器结构恢复接近自然图像的版本，并通过对抗判别器训练，让净化结果在分布上靠近干净图像。公开 README 把管线概括为 poisoned images 到 Detector，再到 Purifier，再到 vanilla 3DGS training，并报告在 Mip-NeRF360 等数据上显著降低高斯数量和峰值显存。它的重要性在于给 3DGS-as-a-service 场景提供了可插拔防线：不要求改动 densification 或渲染器，而是在服务入口处把投毒输入隔离。
相对 Poison-splat 这类攻击论文，RemedyGS 不是再证明风险，而是把防御接口放在训练数据入口，保护后续原始 3DGS 流程不被异常纹理牵引。
\textbf{方法信息：} 代码：\url{https://github.com/Polly-LYP/RemedyGS}。
\textbf{设计动机与边界：} 相对原始 3DGS 以独立高斯和图像重建误差为主的设定，RemedyGS 把问题落在以下可实现的改动上：RemedyGS框架包含三个核心组件：(1)检测器：利用中毒图像和干净图像的局部纹理特征差异（攻击最大化TV分数引入异常高频噪声），设计四层2D卷积+线性分类头的神经网络进行二分类。选择性处理机制确保仅中毒图像被净化，保留正常用户的图像质量。(2)净化器：从中毒输入中恢复干净图像，减少攻击对3DGS重建的不良影响。使用编码器-解码器架构。(3)对抗训练：引入对抗判别器强制恢复图像的分布与原始干净图像对齐，提升感知质量和重建精度。整体框架无需修改3DGS训练过程本身，可即插即用地保护现有3DGS服务。其设计目的不是单独增加一个模块，而是用上述表示、监督或渲染约束消除原流程在目标任务中的具体失配。提出首个针对3DGS计算成本攻击的黑盒防御框架RemedyGS，包含两个关键组件：检测器（基于卷积神经网络识别中毒输入图像的异常纹理模式）和净化器（恢复原始图像并引入对抗训练确保分布对齐）。在白盒、黑盒和自适应攻击下均实现SOTA防御性能，有效平衡安全性和实用性。适用时还应保留论文给出的边界：检测器对新型攻击策略可能失效；净化过程引入额外计算开销；对抗训练需要额外训练数据和时间。
% METHOD: 2603.26067
\paragraph{R-PGA（可重光照 3DGS 物理对抗伪装）（2026.03）}
R-PGA~\cite{2603_26067} 的锚点是自动驾驶中的物理对抗贴图生成。传统物理伪装常在 CARLA 等简化模拟器里优化纹理，结果容易受视角、天气和光照变化影响；R-PGA 的改变是把模拟器换成可重光照的 3DGS，以更真实的物体外观和可微渲染来优化伪装纹理。

具体做法是先用 3DGS 重建目标车辆或物体，再把标准高斯扩展为带物理解耦属性的表示，例如 albedo、normal、roughness 和 metallic，并用 Disney micro-facet BRDF 做前景物体的重光照渲染。背景部分不强求用同一套高斯精确建模，而是通过预训练图像翻译模型生成与前景光照一致的背景，从而形成混合渲染管线。攻击优化时，它不是只最小化平均检测分数，而是加入 Hard Physical Configuration Mining，主动寻找最不利的视角、光照、天气等物理配置，并压低这些配置上的损失峰值，使对抗纹理落在更平坦的损失区域。公开 README 也确认了这一路线：代码依赖 LBM 与 GIR，先做目标物体重建，再生成物理对抗伪装。与前两篇数字投毒不同，R-PGA 把 3DGS 用作物理世界攻击的高保真可微模拟器，目标不是篡改一个资产，而是把可打印或可映射的伪装带到真实拍摄条件中。
\textbf{实现细节：} PGA 的对抗迷彩在优化时固定环境光照，无法适应部署时从正午到黄昏、从晴天到雾霾的连续光照变化；同时，标准 3DGS 将光照烘焙到视角依赖的 SH 系数中，导致跨视角纹理不一致——特定视角的迷彩呈现不同外观，阻碍收敛到统一的物理纹理。R-PGA（2026.03） 作为 PGA 的期刊扩展版本，从场景仿真和优化目标两个维度系统解决了这些局限。在场景仿真层面，R-PGA 首先对 3DGS 的每个高斯基元 $g_j$ 引入显式的物理渲染（PBR）属性：反照率 $a \in [0,1]^{}3$（本质表面颜色，独立于光照）、金属度 $m \in [0,1]$、法线 $n$ 和粗糙度 $r \in [0,1]$，建立物理解耦表示。渲染时，标准 3DGS 的颜色渲染被替换为基于渲染方程的物理着色模型：出射辐射度 $L_o(\omega_o, x) = \int_\Omega f_r(\omega_i, \omega_o, x)\, L_i(\omega_i, x)\, (\omega_i \cdot n)\, d\omega_i$，其中 $f_r$ 采用 Disney 微面元 BRDF 模型，入射光 $L_i$ 分解为直接光照（基于图像的照明 IBL，用高分辨率环境图表示）和间接光照（球谐函数捕获多次反弹，配合可见性项处理遮挡），法线 $n$ 从高斯协方差矩阵的最短轴推导。关键在于：优化反照率 $a$ 而非 SH 系数，使对抗模式反映表面的本质颜色而非瞬时照明效果，从根本上消除了 PGA 中 SH 烘焙光照导致的跨视角纹理不一致问题。然而，对整个场景进行材质分解在开放环境中是不适定的——周边区域视角监督稀疏，导致几何误差、材质/光照分解失败和悬浮伪影相互加剧。
\textbf{实现补充：} 实现上还要核对输入、输出、参数、特征、初始化、训练损失、正则、渲染或推理接口，以及优化中的剪枝、分裂、采样、编码或解码步骤；这些抓手决定该方法到底改了 3DGS 管线的哪一环。
\textbf{方法信息：} 代码：\url{https://github.com/trlou/r-pga}。



\section{数字水印}

水印方向的锚点是“在不明显损害渲染的前提下，把可验证版权消息嵌入 3DGS”。由于 3DGS 是显式高斯集合，水印可以放在高斯参数、球谐颜色、锚点特征，或通过渲染图像再由 2D 解码器读出。不同方法的关键差异在于三个问题：嵌入位置是否容易被压缩或编辑破坏，提取是在 3D 参数域还是 2D 渲染域，是否需要每个场景单独优化。


% METHOD: 2410.23718
\paragraph{GaussianMarker（不确定性感知水印）（2024.10）}
GaussianMarker~\cite{2410_23718} 是本章较早的 3DGS 水印锚点之一。它针对的问题是：直接改高斯参数会在渲染图中产生可见伪影，而完全依赖 2D 输出水印又保护不了被复制的 3DGS 文件。它的改变是先估计各个高斯参数对渲染结果的不确定性，再把版权消息嵌入那些更不敏感、更能容忍扰动的位置。

具体做法是把训练好的 3DGS 参数看作后验分布附近的解，用 Laplace 近似和 Hessian 信息估计参数不确定性。高不确定性参数意味着在小范围扰动下渲染质量更不容易下降，因此水印扰动优先加到这些参数上。训练目标同时包含渲染保真项和消息解码项：一方面让水印后的模型渲染接近原模型，另一方面让解码器能从 3D 高斯参数或 2D 渲染图像中恢复消息。相对普通点云或网格水印，GaussianMarker 的核心不是几何坐标扰动，而是利用 3DGS 参数到图像的可微关系选择可嵌入空间；相对后来的泛化方法，它仍偏逐模型优化，但建立了“3D 参数域加 2D 渲染域双提取”的基本水印框架。
\textbf{实现细节：} 与 3D-GSW 通过 FGD 和梯度掩码来控制水印扰动幅度不同，GaussianMarker（2024.10） 从根本性的新角度出发——利用贝叶斯不确定性来指导水印嵌入，而非在微调过程中通过工程化的掩码约束来限制扰动。3D-GSW 的 FGD 移除低贡献高斯并分裂高频区域高斯，然后对剩余高斯施加基于参数幅度的指数衰减掩码，本质上是一种"全局筛选+逐参数加权"的策略。GaussianMarker 则提出了一种更具理论基础的方案：直接估计 3DGS 模型参数的后验分布，量化每个参数对渲染结果的不确定性贡献，然后仅在高不确定性参数中添加扰动来承载水印信号。
\textbf{方法信息：} NVIDIA 参与了这篇论文；引用很多，Semantic Scholar 约 29 次。



% METHOD: 2411.19895
\paragraph{GuardSplat（CLIP 引导的高效水印）（2024.11）}
GuardSplat~\cite{2411_19895} 同样追求高效水印，但它把重点放在消息解码器的训练方式和球谐特征嵌入上。它的锚点是 3DGS 的 SH 颜色表示：SH 系数控制视角相关外观，通常不改变高斯几何结构；GuardSplat 的改动是把消息嵌入 SH 特征，并用 CLIP 的语义对齐能力帮助训练通用消息解码器，减少逐场景优化成本。

具体做法先用 CLIP-guided message decoupling 训练解码器：二进制消息被映射到 CLIP token 或相关表示，再通过文本/图像对齐能力辅助学习一个更通用的消息空间。嵌入阶段使用 SH-aware 模块，把水印信息调制到每个高斯的球谐颜色特征中，尽量不动位置、尺度和协方差，因此更少破坏几何结构。提取阶段从渲染图像中恢复消息，并在训练中加入 JPEG、裁剪、噪声等 2D 扭曲来增强鲁棒性。项目仓库包含 \texttt{run\_watermark.py}、攻击工具和水印评估脚本，说明它实际面向的是端到端嵌入、攻击模拟和提取评估流程。相对 GS-Marker，GuardSplat 更强调利用预训练视觉语言表示降低解码器训练成本；相对 CompMarkGS，它不专门针对压缩锚点特征，而是把 SH 作为更直接的水印载体。
\textbf{实现细节：} GuardSplat（2024.11） 针对 3DGS 水印在安全性、容量、不可见性和训练效率四方面的平衡挑战提出高效框架。传统方法面临的核心矛盾是：增强水印鲁棒性通常需要更大的参数扰动，但这会降低渲染保真度。与锚点 3D-GSW 直接微调全部高斯参数并用 HiDDeN 架构从零训练消息解码器不同，GuardSplat 在解码器训练范式、水印嵌入策略和鲁棒性增强三个维度上均做出了根本性改进。
\textbf{方法信息：} 代码：\url{https://github.com/narcissusex/guardsplat}。
\textbf{设计动机与边界：} 相对原始 3DGS 以独立高斯和图像重建误差为主的设定，GuardSplat 把问题落在以下可实现的改动上：GuardSplat框架包含三个核心模块：(1)CLIP引导的消息解耦优化——将消息编码器和解码器解耦，利用CLIP的文本-图像对齐能力，将二进制消息映射为CLIP token，通过CLIP文本编码器和MLP解码器训练通用消息解码器，避免逐场景优化。(2)SH-aware消息嵌入——将水印信息直接嵌入每个3D高斯的球谐（SH）特征中，保持原始3D结构不变，使水印难以从模型文件中移除。使用可学习嵌入向量调制SH系数。(3)抗扭曲消息提取——对渲染的2D图像应用各种扭曲（JPEG压缩、裁剪、噪声等）训练鲁棒的消息提取器。其设计目的不是单独增加一个模块，而是用上述表示、监督或渲染约束消除原流程在目标任务中的具体失配。提出首个专为3DGS设计的高效水印框架，通过CLIP引导的解耦优化训练通用消息解码器（非逐场景），SH-aware嵌入保持原始3D结构，抗扭曲提取保证鲁棒性，在容量、安全性、不可见性和效率上达到最优平衡。适用时还应保留论文给出的边界：嵌入容量与不可见性之间存在权衡；对极端攻击（如3D模型直接修改）的鲁棒性有限；需要针对不同3DGS变体适配。
% METHOD: 2503.12836
\paragraph{CompMarkGS（面向压缩的锚点特征水印）（2025.03）}
CompMarkGS~\cite{2503_12836} 的锚点不再是原始 3DGS，而是 Scaffold-GS 一类锚点基表示。它发现已有水印方法即使在未压缩模型上可提取，一旦经过量化或压缩，直接嵌入高斯属性的消息很容易被破坏。它的改变是把水印放进锚点特征而非单个高斯属性，并在训练中显式模拟压缩失真。

具体做法是选用锚点基 3DGS，把消息嵌入 anchor feature。锚点特征再通过 MLP 间接生成局部高斯属性，这让水印不暴露为某个易被手工清除的颜色或不透明度扰动。为适配压缩，它加入 Quantization Distortion Layer，在优化时向特征注入近似量化噪声，让解码器学会从压缩后的表示中恢复消息。渲染质量方面，它使用频率感知锚点增长策略，把更多表示能力放在高频区域，并用 HSV 相关损失压低颜色伪影。项目 README 也说明需要结合 ContextGS 等压缩流程做 compressed extraction。它相对 GaussianMarker 的主要变化是从“不确定参数中嵌入”转向“锚点特征中嵌入并对量化训练”，因此更贴近压缩 3DGS 资产会被实际分发的场景。
\textbf{实现细节：} 在锚点 GaussianSeal 将水印从后处理推进到生成过程嵌入之后，CompMarkGS（2025.03） 将 3DGS 水印推向另一个关键场景——模型压缩。3DGS 模型体积庞大，需要压缩存储和传输。但现有水印方法（包括前述 3D-GSW、WATER-GS、GuardSplat 以及 GaussianSeal）均将水印嵌入高斯属性或网络参数中，一旦经历量化压缩——连续的 SH 系数和几何参数被截断为低精度离散值——水印引入的微小扰动即被抹除。CompMarkGS 首次将压缩鲁棒性作为 3DGS 水印的核心设计目标，使水印在经历量化压缩后仍保持接近无损的比特准确率。CompMarkGS 选择锚点 3DGS （如 Scaffold-GS）作为水印载体，将水印嵌入锚点特征 $f \in \mathbb{R}^d$ 而非直接操作高斯属性。这一选择具有双重优势。安全性上，锚点特征通过 MLP 间接决定高斯属性，攻击者难以通过分析点云或高斯属性直接检测水印的存在。质量上，锚点特征对渲染质量的影响是间接的（经由 MLP 映射），不像位置、缩放、偏移等属性那样直接控制高斯行为，因此修改锚点特征对渲染质量的扰动更可控。具体地，引入与锚点特征同维的可学习水印嵌入特征 $f' \in \mathbb{R}^d$，通过 $f^w = f + \tanh(f')$ 生成含水印的锚点特征。$\tanh$ 函数将扰动限制在 $[-1, 1]$ 的对称范围内。这一设计源自对锚点属性分布的深入分析：训练良好的锚点属性服从正态分布，无界扰动会增大方差导致梯度不稳定和收敛缓慢；而 sigmoid 的 $[0, 1]$ 范围会引入正向偏置；仅 $\tanh$ 的对称有界范围能保证稳定收敛和无偏嵌入。水印的解码与渲染解耦：对渲染图像施加 DWT，将低频 LL 子带输入预训练的 HiDDeN 解码器提取 $L$ 比特消息，水印消息损失为标准二值交叉熵 $\mathcal{L}_{\mathrm{msg}}$。核心创新是量化失真层（QDL，Quantization Distortion Layer）——训练时主动注入与实际压缩算法一致的量化噪声：$\tilde{f}^w_i = f^w_i + \mathcal{U}(-\frac{1}{2}, \frac{1}{2}) \cdot q_i$，其中量化步长 $q_i = Q_0 \cdot (1 + \tanh(r_i))$，$r_i = \mathrm{MLP}_q(f^w_i)$ 通过小型 MLP 根据含水印特征自适应调整，将步长限制在 $[0, 2Q_0]$ 范围内。
\textbf{方法信息：} 代码：\url{https://github.com/kuai-lab/iclr26_compmarkgs}。
\textbf{设计动机与边界：} 相对原始 3DGS 以独立高斯和图像重建误差为主的设定，CompMarkGS 把问题落在以下可实现的改动上：CompMarkGS包含四个核心组件：(1) 锚点基3DGS水印：选择Scaffold-GS架构，水印嵌入到锚点特征f中（f\^{}w = f + tanh(f')），而非直接修改高斯属性。锚点特征通过MLP间接影响渲染质量，攻击者难以直接操纵水印。tanh函数将嵌入特征限制在[-1,1]保证稳定收敛。(2) 量化失真层(QDL)：借鉴可微分失真层思想，在训练过程中模拟量化压缩引入的噪声，使水印学习对量化噪声的鲁棒性。这是3DGS水印研究中首次将可微分失真层应用于模型压缩问题。(3) 频率感知锚点增长策略：识别渲染图像中的高频区域，映射到高斯坐标后选择性应用锚点增长，提升渲染质量。(4) HSV损失：减少颜色伪影，进一步提升渲染质量。其设计目的不是单独增加一个模块，而是用上述表示、监督或渲染约束消除原流程在目标任务中的具体失配。提出CompMarkGS，首个对压缩鲁棒的3DGS水印方法。三大创新：(1) 利用锚点基3DGS将水印嵌入锚点特征（而非直接修改高斯属性），增强安全性和渲染质量；(2) 量化失真层(QDL)在训练时注入量化噪声，使水印在量化压缩后仍保持完整；(3) 频率感知锚点增长策略识别高频区域的高斯进行选择性增长，配合HSV损失减少颜色伪影。适用时还应保留论文给出的边界：仅针对量化基压缩方法的鲁棒性验证；水印容量受限于锚点特征维度；对其他类型攻击（如裁剪、几何变换）的鲁棒性未充分评估。
% METHOD: 2503.18718
\paragraph{GS-Marker（单次前向的泛化水印）（2025.03）}
GS-Marker~\cite{2503_18718} 关注的是水印效率和泛化性。GaussianMarker 这类方法通常需要对每个模型、每条消息做多轮优化；GS-Marker 的改变是把水印嵌入做成单次前向传播，由 3D 编码器把二进制消息写入已优化的 3DGS，再由 2D 解码器从任意视角渲染图像中读出。

具体做法包含三块。第一，3D 编码器类似 PointNet 式地处理无序高斯属性，并把消息嵌入位置、颜色、不透明度、尺度或方向等属性。第二，Adaptive Marker Control 为每类属性引入可学习扰动幅度，解决已训练 3DGS 位于渲染损失局部极小值的问题：扰动太小，消息写不进去；扰动太大，渲染失真。AMC 让模型按属性自适应调节扰动强度。第三，训练中加入畸变层和 2D 解码器，使消息能从渲染图像而不是只从模型文件中恢复。它相对 GaussianMarker 的核心差异是从逐场景优化转向学习到一个可泛化嵌入器，代价是每个高斯都可能参与消息承载，带来一定冗余和容量限制。
\textbf{实现细节：} 与锚点 BitC-3DGS （通过位压缩令牌化在逐场景训练框架内扩展比特容量）不同，可泛化水印的核心挑战是范式转换——从"每个场景都需要一次昂贵的微调"推进到"训练一次、无限场景复用"。GS-Marker（2025.03） 针对逐场景迭代优化的低效和泛化性不足提出单次前馈水印方法。现有 3DGS 水印方法需要对每个场景单独优化，耗时且无法泛化到新场景。GS-Marker 的框架包含三个关键模块：（1）基于 PointNet++ 的 3D 嵌入网络，将归一化的高斯属性与水印消息拼接后，通过级联的局部区域聚合块预测残差调整量，适配 3DGS 的无序点集结构。具体地，嵌入网络输出残差后先投影至 $[-1, 1]$ 范围，再拆分为各高斯属性（位置、缩放、旋转、颜色、不透明度）的独立调整量，分别乘以各属性的缩放因子 $s_*$（位置 $s_M$、缩放 $s_S$、旋转 $s_R$、颜色 $s_C$、不透明度 $s_A$），实现对不同敏感度属性的差异化调整——位置属性需要极小残差而颜色属性可容忍更大修改；（2）失真层，同时模拟 3D 域（dropout、裁剪）和 2D 域（JPEG 压缩、裁剪、旋转）的降质，确保端到端可微训练；（3）基于 HiDDeN 架构的 2D 解码器，从渲染图像中恢复嵌入消息。解码器采用两阶段训练策略：先在输入 3DGS 渲染图像上预训练 HiDDeN 模型以获得良好初始化，再与 3D 编码器端到端联合微调以优化整体性能。其关键创新是自适应标记控制（Adaptive Marker Control, AMC）机制——解决 3DGS 水印优化中的局部最优困境。设计动机是：预训练 3DGS 已经处于渲染损失的局部最优点（梯度极小），直接叠加水印目标会导致水印梯度主导优化方向，破坏渲染保真度或陷入不稳定收敛。AMC 通过在初始 3DGS 上施加受控扰动，将参数推离渲染损失的平坦区域，使渲染保真度和水印目标从优化起始就提供有意义的梯度。
\textbf{方法信息：} Microsoft 参与了这篇论文。
\textbf{设计动机与边界：} 相对原始 3DGS 以独立高斯和图像重建误差为主的设定，GS-Marker 把问题落在以下可实现的改动上：GS-Marker框架包含三个核心模块：(1)3D编码器：使用类PointNet的嵌入网络将二进制消息嵌入3DGS。关键创新是自适应标记控制(AMC)机制，对每个3DGS属性添加可学习尺度的噪声扰动，使其逃离渲染损失的局部极小值，改善训练稳定性和收敛性。还设计了特殊的W-Markers高斯溅射引入随机纹理以增强嵌入效果。(2)畸变层：增强对各种畸变（如压缩、裁剪、噪声等）的鲁棒性。(3)2D解码器：从渲染图像中提取水印消息。联合优化水印解码准确性和渲染保真度，支持任意视角的自由渲染。其设计目的不是单独增加一个模块，而是用上述表示、监督或渲染约束消除原流程在目标任务中的具体失配。提出首个单次前向传播的3DGS水印框架，通过自适应标记控制(AMC)机制解决已优化3DGS的局部极小值问题，实现泛化的消息嵌入和从任意视角渲染图像中提取水印。适用时还应保留论文给出的边界：水印容量受限于高斯基元数量；对极端几何畸变的鲁棒性有待提升；3D编码器对高斯基元分布敏感。

\section{隐写与信息隐藏}

隐写与水印的区别在于目标更强：水印通常只需要证明所有权或恢复短消息，隐写则希望把图像、比特流甚至另一个 3D 场景隐藏在 3DGS 中，同时让外部观察者看不出容器被修改。3DGS 的显式性让隐写更难，因为攻击者可能直接拿到点云文件和高斯属性；因此这一路线的关键是把外观重建和秘密恢复解耦，让同一组高斯属性在普通渲染时像原场景，在秘密解码器下又能恢复隐藏内容。
% METHOD: 2405.15118
\paragraph{GS-Hider（耦合安全特征隐写）（2024.05）}
GS-Hider~\cite{2405_15118} 是 3DGS 隐写路线的重要早期锚点。它从 NeRF 隐写出发，但指出 3DGS 的显式点云文件更透明，不能简单把秘密藏在一个黑盒网络权重里。它的改变是用 coupled secured feature attribute 替代原始 SH 系数，让同一组高斯特征同时承载原始场景外观和隐藏消息，再通过两个不同解码器分别读出。

具体做法是保留高斯的位置、协方差和不透明度等结构属性，把颜色相关的 SH 表达换成维度更高的安全耦合特征。普通渲染时，场景解码器把这些特征解码为 RGB 外观，保证渲染图接近原始 3DGS；秘密提取时，消息解码器从同一特征中恢复隐藏内容，隐藏内容可以是另一个 3D 场景、2D 图像或比特流。训练损失由三部分组成：渲染损失保证容器场景保真，消息恢复损失保证秘密可读，安全或不可检测损失约束耦合特征的统计分布，使其不容易被 steganalysis 发现。项目 README 和代码入口显示它修改了 CUDA rasterizer 以渲染任意维特征图，这与论文方法一致。相对后来的 IBRSteG，GS-Hider 仍更偏每个场景/消息对的优化，但它确立了“替换颜色特征并双解码”的基本范式。
\textbf{实现细节：} GS-Hider（2024.05） 首次提出 3DGS 隐写框架，其核心创新是"耦合安全特征属性"的设计。本方法修改了 3DGS 的颜色表示与渲染管线——用高维耦合特征 $f_i$ 替代原始 SH 系数 $h_i$ 作为可学习颜色属性，同时将 RGB 光栅化替换为耦合特征光栅化——高斯的几何属性（位置 $\mu_i$、旋转 $q_i$、缩放 $s_i$、不透明度 $\alpha_i$）保持不变。从基础隐写框架 GS-Hider 开始，后续工作从生成式隐写、知识蒸馏到频域加密和密钥安全等多个方向进行了拓展。隐写与信息隐藏方法。本节讨论三条技术路线：（1）耦合特征与对比学习隐写；（2）频域加密与双流隐写；（3）密钥安全与混合解耦隐写。
\textbf{设计动机与边界：} 相对原始 3DGS 以独立高斯和图像重建误差为主的设定，GS-Hider 把问题落在以下可实现的改动上：GS-Hider针对3DGS的特殊性（显式表示、透明属性、实时渲染）设计了专门的隐写方案。核心方法是设计一个耦合安全特征属性（coupled secured feature attribute）替代原始3DGS的球谐系数SH。该属性同时编码原始场景外观和隐藏消息，通过两个独立的解码器分别恢复：1）场景解码器：从耦合特征解码出原始RGB渲染图像；2）消息解码器：从耦合特征解码出隐藏的消息（可以是3D场景、图像或比特流）。训练时同时优化三个损失：渲染损失（保证渲染质量）、消息提取损失（保证隐藏信息可恢复）和安全性损失（防止嵌入特征被检测到）。任务设置上，与NeRF隐写不同，GS-Hider假设攻击者可能获取完整的3DGS点云文件，因此安全性要求更高。支持三种隐藏消息类型：3D场景加密、2D图像水印、比特流嵌入。其设计目的不是单独增加一个模块，而是用上述表示、监督或渲染约束消除原流程在目标任务中的具体失配。提出首个面向3DGS的隐写框架GS-Hider，可将3D场景和图像以不可见方式嵌入原始GS点云文件中并准确提取。设计耦合安全特征属性替代原始球谐系数，通过场景解码器和消息解码器解耦原始RGB场景和隐藏消息。支持多模态消息隐藏，在保证渲染质量的同时具备高安全性、鲁棒性、容量和灵活性。适用时还应保留论文给出的边界：嵌入大容量消息时渲染质量略有下降；对极端压缩攻击的鲁棒性有待提升；消息提取需要获取完整点云文件。
\textbf{方法信息：} 代码：\url{https://github.com/xuanyuzhang21/gs-hider}；引用很多，Semantic Scholar 约 34 次。



% METHOD: 2501.03605
\paragraph{ConcealGS（知识蒸馏与梯度引导隐式信息）（2025.01）}
ConcealGS~\cite{2501_03605} 沿着 GS-Hider 的隐写目标继续推进，但它更强调在不牺牲重建质量的情况下把版权或隐式信息稳定嵌入可学习参数。它的锚点是“教师 3DGS 提供原场景渲染，学生 3DGS 同时学习渲染和秘密信息”；改变在于引入知识蒸馏和梯度引导，缓解秘密恢复与渲染保真之间的冲突。

具体做法是先训练或使用一个原始 3DGS 作为教师模型，再训练带隐式信息的学生模型。普通视角上，学生渲染要通过蒸馏损失接近教师渲染，避免因为隐藏信息造成可见失真；检查视角或指定解码条件下，解码器需要从学生渲染或参数中恢复版权标识等隐式信息。为了避免一个目标压过另一个目标，ConcealGS 使用梯度优化策略动态调节解码器各层或各损失的梯度权重，让渲染质量和信息恢复同时收敛。项目 README 将其定位为在 3DGS 中 conceal implicit information，并给出训练、渲染、环境文件，说明它是完整公开实现。相对 GS-Hider 的耦合特征替换，ConcealGS 更像是在教师约束下训练一个带秘密通道的学生高斯场。
\textbf{实现细节：} GaussianStego 面向生成式 3DGS 的生成过程嵌入，而 ConcealGS 则从知识蒸馏的角度提供了一种互补的后处理隐写思路。ConcealGS（2025.01） 的核心区别在于操作时机——GaussianStego 在 3DGS 生成模型的前馈过程中嵌入隐藏信息（通过修改生成网络权重间接影响输出高斯），而 ConcealGS 直接对已训练好的 3DGS 资产进行后处理嵌入，无需访问原始训练数据或生成管线。这样，ConcealGS 适用于更广泛的场景：任何已有的 3DGS 资产都可以直接添加隐式信息，而不依赖特定的生成架构。
\textbf{方法信息：} 代码：\url{https://github.com/zxk1212/concealgs}。
\textbf{设计动机与边界：} 相对原始 3DGS 以独立高斯和图像重建误差为主的设定，ConcealGS 把问题落在以下可实现的改动上：ConcealGS的核心包含三个组件：（1）知识蒸馏策略：使用预训练的3DGS教师模型指导学生模型（含隐式信息）训练，通过渲染一致性损失确保渲染质量不受影响。（2）通用解码器：从检查视角渲染中解码隐式信息，在正常视角上作为恒等网络训练（确保渲染质量不下降），在检查视角上恢复嵌入的信息。使用对比损失确保信息恢复的准确性。（3）梯度引导优化策略：动态调整解码器各层的梯度权重，平衡渲染质量和信息恢复两个目标。处理方案考虑了信息嵌入、覆盖一致性，以及重训练后的信息保留。其设计目的不是单独增加一个模块，而是用上述表示、监督或渲染约束消除原流程在目标任务中的具体失配。提出首个面向3DGS的隐写方法ConcealGS，通过知识蒸馏和梯度优化策略在3DGS可学习参数中嵌入隐式信息。克服了NeRF方法的局限性，增强了隐式信息的鲁棒性和3D重建质量。几乎不影响渲染质量。适用时还应保留论文给出的边界：嵌入容量有限；对恶意攻击（如3DGS参数大幅修改）的鲁棒性有待验证；仅适用于3DGS格式。
% METHOD: 2606.30024
\paragraph{IBRSteG（可泛化秘密场景嵌入）（2026.06）}
IBRSteG~\cite{2606_30024} 是隐写路线中更接近新锚点的方法，因为它把 3DGS 隐写从“每对场景单独优化”改成“前馈可泛化嵌入”。GS-Hider 和 ConcealGS 都要围绕具体 cover scene 与 secret message 做较强优化；IBRSteG 的改变是学习一个 scene-agnostic embedding function，使未见过的 cover scene 和 secret scene 也能被快速处理。

具体流程是先用冻结的 GPS-Gaussian+ 类重建骨干，从 cover 和 secret 的双视图图像得到两组 3DGS。然后通过映射函数把无序高斯属性转成结构化 Gaussian Attribute Maps，例如深度、颜色、旋转和不透明度图；这样 3D 高斯隐写被转换到更适合 2D 网络学习的属性图域。核心模块 Gaussian Attributes Steganographer 在 GAM 域中把 secret scene 的属性嵌入 cover scene，生成 stego Gaussian scene；提取时再从 stego 表示恢复 secret。公开 README 明确说明它隐藏的是完整 secret 3D Gaussian scene，且不需要 per-scene optimization，推理约为秒级，并提供 checkpoint 与训练测试脚本。相对 GS-Hider 的耦合特征，IBRSteG 最大的技术路线变化是引入结构化属性图，把无序高斯集合变成可被前馈网络处理的隐写载体。
\textbf{实现细节：} IBRSteG（2026.06） 是 3DGS 隐写路线中的新增代表。相较于 3DGS，主要改变是：前馈 3D Gaussian embedding 与 Gaussian Attributes Steganography。3DGS 的显式高斯表示和可微渲染能力保持为基础。其设计动机是：将 3D steganography 表述为前馈 3D Gaussian embedding，缓解既有方法与特定场景刚性耦合的问题。具体而言，通过 GAS 等模块把 secret scene 的高斯属性嵌入 cover/stego Gaussian scene，在保持视觉不可察觉的同时支持跨场景泛化。它给后续方法留下的比较点是：如何利用 Gaussian 属性承载隐写信息和秘密三维场景。边界在于：安全性依赖 threat model；压缩、编辑或重训练可能破坏隐藏信息。
\textbf{实现补充：} 实现上还要核对输入数据、输出指标、基线配置、评测协议、系统接口和适用边界；若调用 3DGS 渲染或重建结果，也应说明相机参数、训练设置、优化流程和指标计算如何对齐。
\textbf{技术实现补充：} IBRSteG。IBRSteG（2026.06） 可作为"Steganography for Gaussian scenes"中的新锚点候选。相较于 3DGS，主要改变是：前馈 3D Gaussian embedding 与 Gaussian Attributes Steganography。3DGS 的显式高斯表示和可微渲染能力保持为基础。其设计动机是：将 3D steganography 表述为前馈 3D Gaussian embedding，缓解既有方法与特定场景刚性耦合的问题。具体而言，通过 GAS 等模块把 secret scene 的高斯属性嵌入 cover/stego Gaussian scene，在保持视觉不可察觉的同时支持跨场景泛化。这一锚点判断来自当前分类中的机制关系：候选锚点：安全化，且相关性 0.95；后续方法可围绕 利用 Gaussian 属性承载隐写信息和秘密三维场景。继续比较和扩展。适用边界是：安全性依赖 threat model；压缩、编辑或重训练可能破坏隐藏信息。
\textbf{方法信息：} 代码：\url{https://github.com/lingxiang2023/ibrsteg}。
\begin{takeaway}
本章可以用“信息写入位置”和“威胁模型”来串起来。水印方法从 GaussianMarker 的不确定参数嵌入，发展到 CompMarkGS 的锚点特征抗压缩、GS-Marker 的单次前向泛化、GuardSplat 的 SH 与 CLIP 引导；隐写方法从 GS-Hider 的耦合安全特征，发展到 ConcealGS 的教师蒸馏和 IBRSteG 的前馈秘密场景嵌入。攻击方向则说明 3DGS 的显式性和 densification 既是优势也是风险：GaussTrap 与 StealthAttack 利用视角和密度结构植入幻觉，Poison-splat 利用自适应密度控制放大计算成本，R-PGA 利用高保真可微重光照做物理攻击，而 RemedyGS 把防线前置到输入检测和净化。读者判断这些方法时，应把公司参与、代码、引用和仓库活跃度当作可复核线索，而不是替代方法机制本身。
\end{takeaway}
\textbf{设计动机与边界：} 相对原始 3DGS 以独立高斯和图像重建误差为主的设定，IBRSteG 把问题落在以下可实现的改动上：通过 GAS 等模块把 secret scene 的高斯属性嵌入 cover/stego Gaussian scene，在保持视觉不可察觉的同时支持跨场景泛化。其设计目的不是单独增加一个模块，而是用上述表示、监督或渲染约束消除原流程在目标任务中的具体失配。将 3D steganography 表述为前馈 3D Gaussian embedding，缓解既有方法与特定场景刚性耦合的问题。适用时还应保留论文给出的边界：安全性依赖 threat model；压缩、编辑或重训练可能破坏隐藏信息。
